\section{Pregunta 2. Cifrado autenticado con AES-GCM}
\label{sec:pregunta2}

\begin{enunciado}{(2,5 puntos)}
Utilizando Openssl cifre un fichero usando el algoritmo autenticado AES-GCM
\end{enunciado}

\subsection{Análisis}

El objetivo es proteger un fichero mediante AES (Advanced Encryption Standard) en modo GCM (Galois/Counter Mode), proporcionando confidencialidad, integridad y autenticidad con una misma construcción. Se empleará AES-128-GCM con una clave simétrica $K$ de 16 bytes (128 bits), manteniendo la longitud utilizada en la pregunta anterior y siguiendo la recomendación de las diapositivas del tema 1.2 \ref{ref:slides-gcm}. Se generará una clave específica para este uso.

GCM combina el modo CTR (Counter Mode) con un mecanismo de autenticación. AES cifra bloques de contador y sus resultados se combinan con el contenido original mediante XOR (Exclusive OR) para obtener el criptograma $C$. La etiqueta $T$ se calcula mediante la función GHASH y AES para detectar alteraciones. Los valores internos se obtienen de $K$, sin añadir una clave auxiliar de HMAC (Hash-based Message Authentication Code). Este funcionamiento se explica en los temas 1.2 y 1.4 \ref{ref:tema-gcm}, \ref{ref:tema-aead}.

\begin{figure}[H]
  \centering
  \resizebox{\linewidth}{!}{% Original vector diagram: topic 1.2, slide 46; topic 1.4, pp. 21--23.
% GCM context shares K and N; GHASH includes its prescribed input encoding.
\begingroup
\setstretch{1}
\hyphenpenalty=10000
\newcommand{\gcmblock}[3]{%
  {\fontsize{14}{16}\selectfont\bfseries\boldmath #1}\\[5pt]
  {\fontsize{10}{12}\selectfont\mdseries #2}\\[4pt]
  {\fontsize{9}{11}\selectfont\mdseries (#3)}%
}
\begin{tikzpicture}[x=1cm,y=1cm,
  gcmnode/.style={proceso,text width=35mm,minimum height=18mm,inner sep=7pt},
  gcmpublic/.style={gcmnode,fill=white,funcionpublica},
  gcmsecret/.style={gcmnode,fill=white,funcionsecreta},
  gcmlength/.style={etiqueta,font=\sffamily\fontsize{9}{11}\selectfont,inner sep=3pt}
]
  % Inputs target the outer boundary: both processing paths share this context.
  \draw[funcioncripto,rounded corners=2pt] (3.6,-2) rectangle (12.5,-17.7);
  \node[font=\sffamily\fontsize{14}{16}\selectfont\bfseries]
    at (8,-2.65) {AES-128-GCM};
  \node[font=\sffamily\fontsize{9}{11}\selectfont,align=center]
    at (8,-3.3) {(shared key $K$; $J_0=N\,\Vert\,0^{31}\,\Vert\,1$)};

  \node[gcmsecret] (key) at (1,0)
    {\gcmblock{$K$}{Secret key}{AES-128}};
  \node[gcmpublic] (nonce) at (9,0)
    {\gcmblock{$N$}{Nonce}{unique under the same key}};
  \draw[flujo] (key.east) -- node[gcmlength,above] {16 bytes (128 bits)} (6,0) -- (6,-2);
  \draw[flujo] (nonce.south) -- node[gcmlength,right] {12 bytes (96 bits)} (9,-2);

  \node[gcmnode,funcioncripto] (stream) at (8,-5)
    {\gcmblock{CTR}{Counter Mode}{AES encrypts counter blocks;\\start at $J_0+1$}};
  \node[gcmnode,funcioncripto] (encrypt) at (8,-8.2)
    {\gcmblock{$\oplus$}{Exclusive OR}{plaintext and counter stream}};
  \node[gcmsecret] (plain) at (1,-8.2)
    {\gcmblock{$M$}{Plaintext file}{no padding required}};
  \node[gcmpublic] (cipher) at (15,-8.2)
    {\gcmblock{$C$}{Ciphertext}{same length as $M$}};

  \draw[flujo] (stream.south) -- node[gcmlength,right] {$n$ bytes ($8n$ bits)} (encrypt.north);
  \draw[flujo] (plain.east) -- node[gcmlength,above] {$n$ bytes ($8n$ bits)} (encrypt.west);
  \draw[flujo] (encrypt.east) -- node[gcmlength,above] {$n$ bytes ($8n$ bits)} (cipher.west);

  \node[gcmpublic] (aad) at (1,-11.5)
    {\gcmblock{AAD}{Additional Authenticated Data}{public metadata; not encrypted}};
  \node[gcmnode,funcioncripto] (hash) at (8,-11.5)
    {\gcmblock{GHASH}{Galois-field hash}{$H=\operatorname{AES}_K(0^{128})$;\\process $A$, $C$ and their lengths}};
  \draw[flujo] (encrypt.south) -- node[gcmlength,right] {$C$: $n$ bytes ($8n$ bits)} (hash.north);
  \draw[flujo] (aad.east) -- node[gcmlength,above] {$A$: $a$ bytes\\($8a$ bits)} (hash.west);

  \node[gcmnode,funcioncripto] (mask) at (6,-16)
    {\gcmblock{AES}{Advanced Encryption Standard}{$E=\operatorname{AES}_K(J_0)$}};
  \node[gcmnode,text width=20mm,funcioncripto] (combine) at (11,-16)
    {\gcmblock{$\oplus$}{Exclusive OR}{$T=G\oplus E$}};
  \node[gcmpublic] (tag) at (11,-20)
    {\gcmblock{$T$}{Authentication tag}{integrity and authenticity}};

  \draw[flujo] (hash.south) -- (8,-14.2)
    -- node[gcmlength,above] {$G$: 16 bytes\\(128 bits)} (11,-14.2) -- (combine.north);
  \draw[flujo] (mask.east) -- node[gcmlength,below] {16 bytes\\(128 bits)} (combine.west);
  \draw[flujo] (combine.south) -- node[gcmlength,right,pos=0.8] {16 bytes\\(128 bits)} (tag.north);

  \leyendaFunciones{(8,-21.8)}
\end{tikzpicture}
\endgroup
}
  \caption{Cifrado y autenticación integrados en AES-GCM. La misma clave y el nonce alimentan el contexto de GCM, que produce el criptograma y la etiqueta. Elaboración propia a partir de los temas 1.2 y 1.4 y de la especificación de GCM (\ref{ref:slides-gcm}, \ref{ref:tema-aead}, \ref{ref:nist-gcm}).}
  \label{fig:p2-aes-gcm}
\end{figure}

Se utilizará un nonce público $N$ de 12 bytes (96 bits), generado aleatoriamente para cada operación, y una etiqueta $T$ de 16 bytes (128 bits). El nonce inicializa el contador y no debe repetirse con la misma clave, ya que su reutilización compromete tanto la confidencialidad como la autenticación. Interviene en el cifrado y no actúa como sal de derivación. GCM no requiere relleno y el criptograma conserva el tamaño del original \ref{ref:nist-gcm}.

El resultado conservará $N$, $C$ y $T$. Los metadatos de una posible cabecera se protegerán como AAD (Additional Authenticated Data): información visible cuya integridad se verifica. Esta capacidad corresponde al cifrado AEAD (Authenticated Encryption with Associated Data) del tema 1.4 \ref{ref:tema-aead}. La recuperación utilizará los mismos parámetros y solo entregará el original si se valida la etiqueta; cualquier alteración de los datos protegidos deberá provocar el rechazo.

Se utilizará la interfaz de cifrado \texttt{EVP} de la biblioteca OpenSSL. Aunque los apuntes incluyen un ejemplo con \texttt{openssl enc}, ese comando no admite GCM según la documentación oficial. La biblioteca permite gestionar el nonce y la etiqueta y validar la autenticación \ref{ref:openssl-gcm}.

\subsection{Implementación y validación}
\pendiente{desarrollar la implementación con OpenSSL e incorporar sus salidas y comprobaciones reales.}
